% Generated draft: requires mathematical and source review.
\documentclass[11pt]{article}
\usepackage[margin=25mm]{geometry}
\usepackage{fontspec}
\setmainfont{Latin Modern Roman}
\usepackage{amsmath,amssymb,mathtools,unicode-math}
\setmathfont{Latin Modern Math}
\usepackage{enumitem,needspace,xurl,hyperref}
\setlength{\parindent}{0pt}
\setlength{\parskip}{5pt}
\setlength{\emergencystretch}{3em}
\newcommand{\N}{\mathbb{N}}
\newcommand{\Z}{\mathbb{Z}}
\newcommand{\Q}{\mathbb{Q}}
\newcommand{\R}{\mathbb{R}}
\newcommand{\C}{\mathbb{C}}
\newcommand{\F}{\mathbb{F}}
\newcommand{\A}{\mathbb{A}}
\newcommand{\PP}{\mathbb{P}}
\newcommand{\E}{\mathbb{E}}
\newcommand{\eps}{\varepsilon}
\newcommand{\dd}{\,\mathrm{d}}
\begin{document}
\section*{UnsolvedMath ID: 2303035}
\textbf{Source problem code:} catalogue entry
\subsection{Definitions and mathematical statement}
For \(n\ge3\), write \(\lambda\) for Lebesgue measure on \(\R^n\), and define the origin-centered annulus
\[
A(r_1,r_2)=\{y\in\R^n:r_1<\|y\|<r_2\},\qquad r_1<r_2.
\]
The source calls its closure a closed annulus. With a same-dimensional sphere convention, let \(\sigma\) be surface measure on \(S^{n-1}\subset\R^n\). For measurable \(E\) with \(0<\lambda(E)<\infty\), put
\[
U_E(x)=\frac1{\lambda(E)}\int_E\|x-y\|^{2-n}\,\dd\lambda(y),\qquad
U_S(x)=\frac1{\sigma(S^{n-1})}\int_{S^{n-1}}\|x-y\|^{2-n}\,\dd\sigma(y).
\]
The recorded prior assertion is: if \(\Omega\ne\varnothing\), \(\lambda(\overline\Omega)<\infty\), and \(U_{\overline\Omega}(x)=U_S(x)\) for every \(x\notin\overline\Omega\), then \(\overline\Omega\) is a closed annulus. Armitage and Goldstein ask whether replacing \(\overline\Omega\) by \(\Omega\) throughout yields the conclusion that \(\Omega\) itself is an annulus. Thus the proposed replacement tests all \(x\notin\Omega\), not merely those outside its closure.

Several qualifications are unresolved. Nonemptiness alone does not guarantee positive measure, and measurability of \(\Omega\) is not stated. The source literally writes the comparison sphere in \(\R^m\), while the kernel variable lies in \(\R^n\); \(m=n\) or an embedding is needed. Most materially, the excerpt calls \(\Omega\) an arbitrary subset, without openness or another regularity hypothesis. Adding isolated null-measure points preserves volume potentials and removes points from the exterior test domain, while changing the exact geometry of the set. Hence the asserted rigidity and its proposed extension require geometric regularity or an explicit almost-everywhere set equivalence convention; these cannot be silently supplied. The permitted radii and treatment of a zero or negative inner radius also need clarification if balls or punctured balls are relevant.

The supplied OPEN-TRIAGE summary reports open status in the 2018 edition and no recent resolution located. The prior assertion is retained as source data with missing hypotheses, not certified as a theorem for arbitrary subsets.
\subsection{Short English statement}
Does the Newtonian-potential characterization stated for closures characterize the original set as an annulus as well? The excerpt has a sphere-dimension mismatch and lacks positive-measure, measurability, and geometric regularity hypotheses needed for exact-set rigidity.
\subsection{Sources}
\begin{itemize}
\item Local UnsolvedMath registry, permanent ID \(2303035\), exact source code \texttt{catalogue entry}, ``Research Problems in Function Theory --- Problem 3.35.''
\item Hayman, \emph{Research Problems in Function Theory} (2018), Problem 3.35, attributed to D. H. Armitage and M. Goldstein; \url{https://arxiv.org/abs/1809.07200}. Source-TeX extraction; recorded access 2026-07-29.
\item The transcription says ``see and'' for related work, with the references missing. No authoritative supporting citation for the stated characterization or current-status resolution is supplied beyond the problem list.
\end{itemize}
\end{document}
